\documentclass[main.tex]{subfiles}


\begin{document}

%from pagenumber to up
% \phantomsection 
% \hypertarget{toc}{}

 

 

% номер секции вручную
\setcounter{section}{35
}
\addtocounter{section}{-1} 

\section{Сила Ампера}


\textbf{Сила Ампера} ($F_\text{А}$ [Н])~--- это сила, действующая на провод с током со стороны магнитного поля:
\begin{equation}
    F_\text{А}=IBl\sin\alpha,
\end{equation}
где $I$~--- сила тока в проводе, $B$~--- индукция магнитного поля, $l$~--- длина провода (находящегося в магнитном поле), $\alpha$~--- угол между направлениями тока~$I$ и поля~$B$.

% дуговой магнит (однородное пурышева 8)

Пусть провод с током помещен между полюсами дугового магнита (рис.\,\ref{fig:p153}).

    \begin{figure}[H]
    \centering
\includegraphics[]{el/39aFA_pic1.pdf}%
\caption{Сила Ампера}
\label{fig:p153}
\end{figure}


% Провод с током~$I$, находящийся в магнитном поле~$B$, созданном дугообразным магнитом, испытывает действие со стороны магнитного поля с силой~$F_\text{А}$.

Поле магнита с индукцией~$B$ (его линии изображены серым цветом) действует на провод с током~$I$ с силой Ампера~$\vec F_\text{А}$. Направление силы Ампера можно определять по следующему правилу.
%
\begin{law}
    \textbf{Правило левой руки.} Если расположить левую руку так, чтобы вытянутые четыре пальца указывали направление тока, а линии поля входили в ладонь, то отведенный перпендикулярно большой палец (лежащий в одной плоскости с остальными пальцами) укажет направление силы Ампера.
\end{law}
%
Рисунок~\,\ref{fig:p154} иллюстрирует применение этого правила.  

    \begin{figure}[H]
    \centering
\includegraphics[]{el/39aFA_pic2.pdf}%
\caption{Правило левой руки}
\label{fig:p154}
\end{figure}





























 
\end{document}